\begin{textsample}{NEG Dim 5 – ba – Score -3.00 – ba/ba\_inf\_8.txt}  \label{ex:f5_neg_004}
7.8.
Transversalidades Fundantes e Transversalidades das Competências da Base Transversalidades fundantes : concepção de infância, cuidado, interação, ludicidade e formação.
Enquanto \textbf{transversalidade}, as políticas curriculares e as atitudes e práticas advindas da sua apropriação deverão deixar evidentes de que concepção de infância partem suas concepções curriculares e como cultivam essa concepção nos seus atos de currículo.
Como, por exemplo, superam as concepções pedagógicas adultocêntricas de criança e/ou reduzidas ao etapismo psicológico, que imaginou sempre a criança como ser imaturo e insuficiente, tendo o “ amadurecimento ” adulto como meta.
Como apreendem, por outro lado, as crianças na sua inteireza existencial/relacional, social, ética, estética, cultural e política e, em termos curriculares, como, ademais, um ator/atriz curriculante.
Ou seja : que acrescenta ao currículo que experiencia.
Da perspectiva do cuidado como uma \textbf{transversalidade} curricular, mas presente na constituição mesma em toda educação de crianças, deixar claro que o cuidado é parte de todo e qualquer ato de currículo.
O cuidado deve estar no cuidar, na mediação das aprendizagens, das interações cognitivas, afetivas, artísticas e linguísticas.
É assim que essa \textbf{transversalidade} insubstituível qualifica um currículo em Educação Infantil.
No que se refere às interações, é fundamental saber acolher, compreender, orientar e reorientar as experiências curriculares das crianças em toda e qualquer atividade que envolva o currículo e seus atos.
As interações saudáveis e acolhedoras qualificam, por consequência, o processo formativo veiculado pelos atos de currículo.
Interações qualificadas habilitam também a aprendizagem e o desenvolvimento das crianças em toda a sua integralidade.
Até porque, como todo o ser humano, a criança é um ser constituído pela sua individualidade, mas, em muito, pelas interações que a constituem em toda a sua vida de relação com saberes, assim como, e por consequência, em relação às suas aprendizagens.
Compreendida como experiência singular de cada criança, a formação, construída como possibilidade de qualificação das aprendizagens, integra a ideia de que uma aprendizagem boa é aquela que qualifica a vida da própria criança.
Implica, portanto, em boas mediações para se chegar a uma formação qualificada, que tem na sua dinâmica a autoformação, formação consigo mesma como ser reflexivo e crítico, heteroformação, formação com o outro ( pais, professores, colegas, autores, personagens da literatura etc. ), ecoformação ( formação na relação com as experiências com outros seres ), metaformação, reflexão sobre o próprio processo formativo.
Um currículo de Educação Infantil só se legitima se qualificar a formação da criança nos níveis da sua integralidade, envolvendo saberes de campos de conhecimento sistematizados, incluindo saberes e experiências éticas, estéticas, culturais, políticas e espirituais.
Quanto à ludicidade, emerge num currículo de Educação Infantil como uma \textbf{transversalidade} que vai ao encontro da própria forma de configuração das experiências infantis vinculadas às brincadeiras e aos jogos e suas potencialidades formativas.
Um currículo de Educação Infantil deverá ser sempre um currículo brincante para um ser brincante ( CONCEIÇÃO ; MACEDO, 2018 ).
É necessário realçar que afirmar que a criança aprende brincando é algo que ultimamente vem sendo discutido com afinco nos meios educacionais.
Da nossa perspectiva, o brincar é uma característica fundamental do ser humano, do qual a criança depende muito para se desenvolver.
Portanto, para crescer, ela deve brincar, bem como para se posicionar frente ao mundo.
Uma preocupação emerge quando a tendência em nossas escolas é de adotar a brincadeira com o caráter de preparo, e não com o caráter recreativo.
Essa realidade faz com que a utilização problematizada da brincadeira na escola se torne premente.
Trazer o lúdico como brincadeira de forma vivenciada e reflexiva de volta para o cotidiano educacional das crianças, salvá -lo de certa prática instrumentalista é um dos aspectos importantes para o qual tentamos chamar a atenção.
Tomar consciência desse processo requer, na verdade, mudanças nas atitudes dos atores e sujeitos do currículo.
Essas mudanças, porém, não acontecem de forma automática.
São necessárias as vivências reflexivas para incorporar a inspiração lúdica vinculada à brincadeira.
Talvez seja necessário reinspirar e reencantar na escola o brincar na educação das crianças.
Com o corpo, com o espaço e os objetos ; com a imaginação, a criatividade, a inteligência ; com a intuição, com as palavras e com o conhecimento, com nós mesmos e com os outros.
Assim, estaremos redescobrindo também essa linguagem para nos comunicar e nos expressar.
Tanto o brincar como o faz de conta podem ser considerados atividades fundamentais para o desenvolvimento da autonomia, uma das finalidades inelimináveis da educação e, a fortiori, da educação das crianças.
Nas brincadeiras, as crianças podem desenvolver algumas capacidades importantes, tais como a atenção, a imitação, a memória, a imaginação.
Amadurecem também algumas capacidades de socialização, por meio da interação e da utilização e experimentação de regras e papéis sociais vivenciados nas brincadeiras.
Essas compreensões já estão presentes nos Referenciais Curriculares da Educação Infantil.
Falta, entretanto, que estas inspirações entrem na vida dos currículos de Educação Infantil de forma tal que representem o início da configuração de uma cultura e de uma práxis educativa que atentem para a multiplicidade das expressões lúdicas dos segmentos infantis ; edifiquem a ludicidade curricular da educação enquanto uma construção identitária.
A escola, enquanto espaço institucionalmente organizado, que contém regras específicas para o seu funcionamento, em geral, impõe à criança e, logicamente, à ludicidade, a sua óptica.
A inserção do processo dicotômico entre o brincar e o trabalho, entre ludicidade e produtividade, frente à fragilidade da significação desta perspectiva de ludicidade, por sua vez, impele a escola a buscar uma “ formação ” puramente cognitiva e utilitarista, marcando a vida da criança por meio da imitação de ações produtivistas, como forma de preparação para uma escolaridade reducionista.
Claramente estabelece uma oposição em relação aos anseios lúdicos infantis.
O mundo produtivo, a cultura ocidental e, consequentemente, a escola vão destituindo as atividades lúdicas de seu valor, considerando -as apenas como uma atividade de descanso em face do excesso de energia.
Um estigma para a vivência do lúdico.
Na escola, a criança manipula os brinquedos que são oferecidos, mas temos dúvidas se ela brinca, porque, muitas vezes, não se reconhece nessa experiência.
O lúdico deixou de ser a totalidade que se constrói para nela habitar.
As práticas da fantasia, por exemplo, que deveriam ser acolhidas, respeitadas e estimuladas antes da entrada em cena do discurso racional, são, muitas vezes, esvaziadas ( CONCEIÇÃO ; MACEDO, 2018 ).
Importante mencionar as “ Dez Competências ” definidas na BNCC : - Valorizar e utilizar os conhecimentos historicamente construídos sobre o mundo físico, social, cultural e digital para entender e explicar a realidade, continuar aprendendo e colaborar para a construção de uma sociedade justa, democrática e inclusiva ; - Exercitar a curiosidade intelectual e recorrer à abordagem própria das ciências, incluindo a investigação, a reflexão, a análise crítica, a imaginação e a criatividade, para investigar causas, elaborar e testar hipóteses, formular e resolver problemas e criar soluções ( inclusive tecnológicas ) com base nos conhecimentos das diferentes áreas ; - Valorizar e fruir as diversas manifestações artísticas e culturais, das locais às mundiais, e também participar de práticas diversificadas da produção artístico-cultural. - Utilizar diferentes linguagens – verbal ( oral ou visual-motora, como Libras, e escrita ), corporal, visual, sonora e digital –, bem como conhecimentos das linguagens artísticas, matemática e científica, para se expressar e partilhar informações, experiências, ideias e sentimentos em diferentes contextos e produzir sentidos que levem ao entendimento mútuo ; - Compreender, utilizar e criar tecnologias digitais de informação e comunicação de forma crítica, significativa, reflexiva e ética nas diversas práticas sociais ( incluindo as escolares ), para se comunicar, acessar e disseminar informações, produzir conhecimentos, resolver problemas e exercer protagonismo e autoria na vida pessoal e coletiva ; - Valorizar a diversidade de saberes e vivências culturais e apropriar -se de conhecimentos e experiências que lhe possibilitem entender as relações próprias do mundo do trabalho e fazer escolhas alinhadas ao exercício da cidadania e ao seu projeto de vida, com liberdade, autonomia, consciência crítica e responsabilidade ; - Argumentar, com base em fatos, dados e informações confiáveis, para formular, negociar e defender ideias, pontos de vista e decisões comuns que respeitem e promovam os direitos humanos, a consciência socioambiental e o consumo responsável em âmbito local, regional e global, com posicionamento ético em relação ao cuidado de si, dos outros e do planeta. - Conhecer -se, apreciar -se e cuidar de sua saúde física e emocional, compreendendo -se na diversidade humana e reconhecendo suas emoções e as dos outros, com autocrítica e capacidade para lidar com elas ; - Exercitar a empatia, o diálogo, a resolução de conflitos e a cooperação, fazendo -se respeitar, promovendo o respeito ao outro e aos direitos humanos, com acolhimento e valorização da diversidade de indivíduos e de grupos sociais, seus saberes, identidades, culturas e potencialidades, sem preconceitos de qualquer natureza ; - Agir pessoal e coletivamente com autonomia, responsabilidade, flexibilidade, resiliência e determinação, tomando decisões com base em princípios éticos, democráticos, inclusivos, sustentáveis e solidários.
Essas competências são aqui apropriadas como \textbf{transversalidades} que podem emergir em todos os argumentos do Currículo do Estado da Bahia.
Ou seja, compreendidas como saberes em uso, nos quais saberes, habilidades e valores compõem as competências transversais a serem desenvolvidas, entretecidas à formação política, ética, estética, cultural, espiritual e religiosa, configuram -se como uma continuidade a ser cultivada pelos atos de currículo durante toda a Educação Básica, respeitando e acolhendo as especificidades de cada etapa.
As \textbf{transversalidades} pontuam e atravessam, como espirais, as etapas formativas, produzindo conexões diversas com as competências, seus componentes e todo e qualquer conteúdo ou atividades vinculadas à formação como um todo.
São, em realidade, macroconcepções que acolhem e também se conectam a todos os atos de currículos e, por consequência, às experiências formativas.
Entretanto, faz -se necessário afirmar, como ressaltamos em nossos argumentos ao longo desse Referencial, que existem outras \textbf{transversalidades} que não implicam tão somente no que a BNCC preconiza, ou mesmo poderá ressignificar o que a BNCC pleiteia como Competência a ser trabalhada, levando em conta a singularidade dos contextos de formação.
O exercício importante é se perguntar : que \textbf{transversalidades} emergem no meu contexto curricular-formacional que a BNCC não pleiteia como \textbf{transversalidade} nas suas competências?
Como podemos trabalhar seu enriquecimento, sua crítica, sua reconfiguração e sua complementação?

% matched lemmas: transversalidade
\end{textsample}
